\scanpage{151}
\mspage{4}
% archive foliation: 150
$\overline{\underline{\mathrm{IV}}}$
\uline{Voisinages tubulaires d'une partie}
modérée $Y$ d'un espace modéré $X$.

C'est par définition un cône strict
autour de $o$ de $X/Y$.

Construction : Soit $f:X\unsure{-|\longrightarrow}\mathbb R^+$
(fonction modérée) telle que
\[
f^{-1}(\{0\})=Y.
\]
On prend $V=f^{-1}([0,\varepsilon])$
$\varepsilon$ petit [et admissible], et
\[
\dot V=f^{-1}(\{\varepsilon\}).
\]

\margin{Choix du vois. tubulaire,
ou du vois. tub. rigidifié, est \uline{anodin}.
\unsure{Noter} que $\dot V$ bien défini en cat.
$\infty$-\unsure{isotopiques}\ldots}

Rigidification du voisinage tubulaire
($=$ du cône). Rigidification
$Y$-admissible : s'insère dans diagramme \inserted{(\unsure{comm})}

\diagram{\begin{tikzcd}[column sep=small,row sep=2em]
V & V-Y=V^*\simeq\dot V\times[0,1[
 \arrow[dl,leftrightarrow] \arrow[dr] & \\
X\simeq \arrow[dr] &
\bigl((X\setminus(\unsure{\dot V|V}\setminus\dot V))\amalg(\dot V\times I\deleted{\unsure{\varepsilon}})\bigr)
& \deleted{\unsure{\amalg}}\underset{\dot V\times\{1\}}{\unsure{\vee|\amalg}}\deleted{\unsure{\amalg}}Y, \\
& Y \arrow[ur,leftrightarrow] &
\end{tikzcd}}

par $\varphi:\dot V\times\{1\}\longrightarrow Y$ appl. modérée
(déterminée de façon unique par la
rigidification).

a) Rigidification $Y$-admissible existe

b) La classe d'homotopie de $\dot V\longrightarrow Y$
bien déterminée

Mais on voudrait plus précis.
